% Part: normal-modal-logic
% Chapter: filtrations
% Section: more-filtrations
%READERNOTE{SOL6-A173: अतिरिक्त अटीमुळे प्राप्यतेच्या जोड्या वाढत नाहीत; त्या कमी होतीलच असे नाही. दोन्ही निस्यंदने समान होऊ शकतात, म्हणून अधिक सूक्ष्मपणाच्या तुलनेत समतेची शक्यता स्पष्ट केली.}

\documentclass[../../../include/open-logic-section]{subfiles}

\begin{document}

\olfileid{nml}{fil}{acc}

\olsection{निस्यंदने आणि प्राप्यतेचे गुणधर्म}

पाहिल्याप्रमाणे, वैश्विक, सर्वत्र उत्तरवर्ती आणि परावर्ती
प्रतिमानांची निस्यंदने अनुक्रमे वैश्विक, सर्वत्र उत्तरवर्ती
आणि परावर्ती असतात. पण सममित किंवा संक्रमक प्रतिमानाचे
प्रत्येक निस्यंदन अनुक्रमे सममित किंवा संक्रमक असेलच असे नाही.
तरी काही बाबतीत हे गुणधर्म टिकतील अशी निस्यंदने
परिभाषित करता येतात. त्यासाठी सर्वांत स्थूल निस्यंदनाच्या
व्याख्येसारखीच पद्धत वापरू, पण~$R^*$ च्या व्याख्येत
अतिरिक्त अटी घालू. $\Gamma$ उप-!!{formula}s खाली बंद
आहे असे समजा. प्रतिमान~$\mModel{M} =\tuple{W, R, V}$
मधील जगांमध्ये $u$, $v$ लागू होणारे
\olref{tab:Cn-filtrations} मधील संबंध~$C_i(u,v)$ पाहा.
या अटींच्या संयोगांवरून $R^*[u][v]$ परिभाषित करता येतो.
उदाहरणार्थ, $R^*[u][v]$ तेव्हा आणि तेव्हाच $C_1(u,v)$
ही अट पूर्ण होते असे ठरवल्यास नेमके सर्वांत स्थूल
निस्यंदन मिळते. $R^*[u][v]$ तेव्हा आणि तेव्हाच
$C_1(u,v)$ आणि $C_2(u, v)$ या दोन्ही अटी पूर्ण होतात
असे ठरवल्यास अधिक सूक्ष्म किंवा तेवढेच सूक्ष्म निस्यंदन मिळते. केवळ $C_1(u,v)$
पेक्षा $C_1(u,v)$ आणि $C_2(u,v)$ या दोन्ही अटी पूर्ण
करणाऱ्या जगांच्या जोड्या कमी किंवा तितक्याच असल्याने हे निस्यंदन
सर्वांत स्थूल निस्यंदनापेक्षा “किमान तितके सूक्ष्म” असते.

\begin{table}[ht]
  \centering
  \begin{tabular}{|ll|}
    \hline
    \multirow{2}{*}{$C_1(u,v)$:}
    \iftag{prvBox}{&
      $\Box!A \in \Gamma$ आणि $\mSat{M}{\Box!A}[u]$ असेल,
      तर $\mSat{M}{!A}[v]$;\iftag{prvDiamond}{ आणि}{}\\}{}
    \iftag{prvDiamond}{&
      $\Diamond!A \in \Gamma$ आणि $\mSat{M}{!A}[v]$ असेल,
      तर $\mSat{M}{\Diamond!A}[u]$; \\}{}
    \hline
    \multirow{2}{*}{$C_2(u,v)$:}
    \iftag{prvBox}{&
      $\Box!A \in \Gamma$ आणि $\mSat{M}{\Box!A}[v]$ असेल,
      तर $\mSat{M}{!A}[u]$;\iftag{prvDiamond}{ आणि}{}\\}{}
    \iftag{prvDiamond}{&
      $\Diamond!A \in \Gamma$ आणि $\mSat{M}{!A}[u]$ असेल,
      तर $\mSat{M}{\Diamond!A}[v]$; \\}{}
    \hline
    \multirow{2}{*}{$C_3(u,v)$:}
    \iftag{prvBox}{&
      $\Box!A \in \Gamma$ आणि $\mSat{M}{\Box!A}[u]$ असेल,
      तर $\mSat{M}{\Box!A}[v]$;\iftag{prvDiamond}{ आणि}{}\\}{}
    \iftag{prvDiamond}{&
      $\Diamond!A \in \Gamma$ आणि $\mSat{M}{\Diamond!A}[v]$ असेल,
      तर $\mSat{M}{\Diamond!A}[u]$; \\}{}
    \hline
    \multirow{2}{*}{$C_4(u,v)$:}
    \iftag{prvBox}{&
      $\Box!A \in \Gamma$ आणि $\mSat{M}{\Box!A}[v]$ असेल,
      तर $\mSat{M}{\Box!A}[u]$;\iftag{prvDiamond}{ आणि}{}\\}{}
    \iftag{prvDiamond}{&
      $\Diamond!A \in \Gamma$ आणि $\mSat{M}{\Diamond!A}[u]$ असेल,
      तर $\mSat{M}{\Diamond!A}[v]$; \\}{}
    \hline
    \end{tabular}
  \caption{निस्यंदने परिभाषित करण्यासाठी संभाव्य जगांवरील अटी.}
  \ollabel{tab:Cn-filtrations}
\end{table}

\begin{thm}\ollabel{thm:more-filtrations}
  $\mModel{M} =\tuple{W,R,V}$ हे प्रतिमान आणि
  $\Gamma$ हा उप-!!{formula}s खाली बंद संच असू द्या.
  $W^*$ आणि $V^*$ यांची व्याख्या
  \olref[fil]{defn:filtration} प्रमाणे करा. मग:
  \begin{enumerate}
  \item $R^*[u][v]$ तेव्हा आणि तेव्हाच $C_1(u, v) \land
    C_2(u,v)$ असे समजा. मग $R^*$ सममित आहे; आणि
    $\mModel{M}$ सममित असेल, तर $\mModel{M^*} =
    \tuple{W^*,R^*,V^*}$ हे निस्यंदन आहे.
  \item $R^*[u][v]$ तेव्हा आणि तेव्हाच $C_1(u, v)
    \land C_3(u,v)$ असे समजा. मग $R^*$ संक्रमक आहे;
    आणि $\mModel{M}$ संक्रमक असेल, तर
    $\mModel{M^*}=\tuple{W^*,R^*,V^*}$ हे निस्यंदन आहे.
  \item $R^*[u][v]$ तेव्हा आणि तेव्हाच $C_1(u, v) \land C_2(u,v)
    \land C_3(u,v) \land C_4(u,v)$ असे समजा. मग $R^*$
    सममित आणि संक्रमक आहे; आणि $\mModel{M}$ सममित
    व संक्रमक असेल, तर $\mModel{M^*}=\tuple{W^*,R^*,V^*}$
    हे निस्यंदन आहे.
  \item $R^*$ ची व्याख्या $R^*[u][v]$ तेव्हा आणि तेव्हाच $C_1(u,
    v) \land C_3(u,v) \land C_4(u,v)$ अशी करा. मग $R^*$
    संक्रमक आणि युक्लिडी आहे; आणि $\mModel{M}$ संक्रमक
    व युक्लिडी असेल, तर $\mModel{M^*}=\tuple{W^*,R^*,V^*}$
    हे निस्यंदन आहे.
  \end{enumerate}
\end{thm}

\begin{proof}
  \begin{enumerate}
    \item $R^*$ सममित आहे हे लगेच दिसते, कारण $C_1(u,v)
      \Leftrightarrow C_2(v,u)$ आणि $C_2(u,v) \Leftrightarrow
      C_1(v,u)$. त्यामुळे $\mModel{M}$ सममित असेल,
      तर $\mModel{M^*}$ हे $\Gamma$ मधून मिळणारे
      निस्यंदन आहे, एवढेच दाखवायचे आहे.
      $C_1(u,v)$ या अटीमुळे
      \iftag{prvBox}{%
        \iftag{prvDiamond}{%
          \olref[fil]{defn:filtration-R2} आणि
          \olref[fil]{defn:filtration-R3} या
          \olref[fil]{defn:filtration} मधील अटी}{%
          \olref[fil]{defn:filtration}\olref[fil]{defn:filtration-R2} ही अट}}{%
        \olref[fil]{defn:filtration}\olref[fil]{defn:filtration-R3} ही अट}
      पूर्ण होते. म्हणून आता फक्त
      \olref[fil]{defn:filtration}\olref[fil]{defn:filtration-R1},
      म्हणजे $Ruv$ वरून $R^*[u][v]$ मिळते हे तपासायचे आहे.

      $Ruv$ असे समजा. $R^*[u][v]$ दाखवण्यासाठी
      $C_1(u,v)$ आणि $C_2(u,v)$ सिद्ध कराव्या लागतील.
      $C_1$ साठी: \iftag{prvBox}{$\Box!A
        \in\Gamma$ आणि $\mSat{M}{\Box!A}[u]$ असेल,
        तर $Ruv$ असल्याने $\mSat{M}{!A}[v]$ सुद्धा.
        \iftag{prvDiamond}{ त्याचप्रमाणे,
        }{}}{}\iftag{prvDiamond}{$\Diamond!A \in \Gamma$
        आणि $\mSat{M}{!A}[v]$ असेल, तर $Ruv$
        असल्याने $\mSat{M}{\Diamond!A}[u]$.}{}
      $C_2$ साठी: \iftag{prvBox}{$\Box!A \in \Gamma$
        आणि $\mSat{M}{\Box!A}[v]$ असेल, तर सममितीमुळे
        $Ruv$ वरून $Rvu$; म्हणून $\mSat{M}{!A}[u]$.
        \iftag{prvDiamond}{ त्याचप्रमाणे,
        }{}}{}\iftag{prvDiamond}{$\Diamond!A \in\Gamma$
        आणि $\mSat{M}{!A}[u]$ असेल, तर
        $\mSat{M}{\Diamond!A}[v]$ (कारण सममितीमुळे $Rvu$).}{}
    \item सराव.
    \item सराव.
    \item सराव.
  \end{enumerate}
\end{proof}

\begin{prob}
  \olref[nml][fil][acc]{thm:more-filtrations} ची सिद्धता पूर्ण करा.
\end{prob}

\end{document}
